\documentclass[12pt]{article}
\usepackage[T1]{fontenc}
\usepackage[utf8]{inputenc}
\usepackage{lmodern}
\usepackage{textcomp}
\usepackage{enumitem}
\usepackage{geometry}
\usepackage{graphicx}
\usepackage{amsmath, amssymb}
\geometry{margin=1in}
\begin{document}
\begin{center}
Astronomy - Graduate Level Assessment 9 \
Total Marks: 40 \
Answer all questions.
\end{center}

\section*{Multiple Choice (10 marks)}
\begin{enumerate}[label=\arabic*.]
\item Mars might be a place for future human explorations. However humans can not breathe on the surface of Mars because the atmosphere consists mostly of ...
\begin{enumerate}[label=(\alph*)]
    \item Nitrogen
    \item Argon
    \item Methane
    \item CO 2
\end{enumerate}
\item What would weigh the most on the moon?
\begin{enumerate}[label=(\alph*)]
    \item A kilogram of feathers
    \item Five pounds of bricks as measured on Earth
    \item Five kilograms of feathers
    \item A kilogram of bricks
\end{enumerate}
\item You've made a scientific theory that there is an attractive force between all objects. When will your theory be proven to be correct?
\begin{enumerate}[label=(\alph*)]
    \item The first time you drop a bowling ball and it falls to the ground proving your hypothesis.
    \item After you've repeated your experiment many times.
    \item You can never prove your theory to be correct only "yet to be proven wrong".
    \item When you and many others have tested the hypothesis.
\end{enumerate}
\item One astronomical unit (AU) is equal to approximately ...
\begin{enumerate}[label=(\alph*)]
    \item 130 million km
    \item 150 million km
    \item 170 million km
    \item 190 million km
\end{enumerate}
\item The visible part of the electromagnetic spectrum is between ...
\begin{enumerate}[label=(\alph*)]
    \item 240 to 680 nm.
    \item 360 to 620 nm.
    \item 380 to 740 nm.
    \item 420 to 810 nm.
\end{enumerate}
\end{enumerate}

\section*{True / False (10 marks)}
\textit{Answer True or False}
\begin{enumerate}[label=\arabic*.]
\setcounter{enumi}{5}
\item True or False: The Chicxulub Crater, located on the Yucatan Peninsula in Mexico, is linked to the mass extinction event that eliminated the dinosaurs 65 million years ago.
\item True or False: The U.S.S.R. was the first country to successfully send landers to the surface of Venus during the Venera program.
\item True or False: Titan is the only moon in the outer solar system known to possess a thick atmosphere primarily composed of hydrocarbons.
\item True or False: The ratio of solar radiation flux on Mercury's surface at perihelion to aphelion is 4:1.
\item True or False: A lunar day, defined as the time from one sunrise to the next, lasts approximately 29 Earth days.
\end{enumerate}

\section*{Long Form Response (20 marks)}
\textit{Answer each question with a clear and complete explanation.}
\begin{enumerate}[label=\arabic*.]
\setcounter{enumi}{10}
\item Discuss the essential properties of our solar system that formation theories must address, and critically analyze why the presence of life on Earth is not a necessary component of these theories.
\item Discuss the methods used to determine the distance of the rock identified as "bigfoot" on Mars, given its size and angular measurement. How do these calculations enhance our understanding of Martian terrain?
\end{enumerate}

\vfill
\noindent\textit{End of Paper}
\end{document}